\documentclass{article}

\usepackage[final]{neurips_2019}

\usepackage[utf8]{inputenc}
\usepackage[T1]{fontenc}
\usepackage{hyperref}
\usepackage{url}
\usepackage{booktabs}
\usepackage{amsfonts}
\usepackage{amsmath}
\usepackage{graphicx}
\usepackage{xcolor}

\setcitestyle{round}

\newcommand{\note}[1]{\textcolor{blue}{{#1}}}

\title{
  State-Dependent Macroeconomic Forecasting \\
  for Equity Returns \\
  \vspace{1em}
  \small{\normalfont Stanford CS229 Project}
}

\author{
  Jagat Brahma \\
  Department of Computer Science \\
  Stanford University \\
  \texttt{jbrahma@stanford.edu}
}

\begin{document}

\maketitle

\begin{abstract}
Predicting the stock market from economic data is notoriously difficult -
the signal is weak and easily drowned out by noise. This project builds a
multi-stage machine learning pipeline called SSRF (State-Dependent Supervised
Screening \& Regularized Factor model) that forecasts S\&P~500 monthly returns
from over 100 macroeconomic indicators. The model achieves~70\% directional
accuracy out-of-sample from 2000--2026, with a Sharpe ratio of~0.51 and strong
statistical significance ($p<0.0001$). An asymmetric leverage strategy
(2.5$\times$ long~/~0.25$\times$ short) further improves the Sharpe to~0.84,
providing meaningful risk-adjusted returns above buy-and-hold.
\end{abstract}


\section{Introduction}

Can economic data predict the stock market? This question has fascinated
researchers for decades \citep{goyal2008comprehensive}. The challenge is
that financial markets are noisy and driven by factors far beyond
government-reported statistics. Yet intuitively, when the economy is strong
(employment is high, factories are busy, consumers are spending), corporate
profits-and therefore stock prices-should rise.

The input to our system is a set of roughly~130 monthly economic statistics
published by the Federal Reserve and other agencies: employment, industrial
production, housing starts, interest rates, consumer sentiment, and inflation.
The output is a prediction of whether the S\&P~500 index (the most widely-
followed U.S. stock market benchmark) will go up or down next month, and by
roughly how much. We designed a four-stage pipeline that filters out noisy
features, extracts hidden economic factors, and adjusts predictions based on
market conditions. We evaluate using ``walk-forward'' backtesting, simulating
real-time performance over the last~26 years with no future data leakage.


\section{Related Work}

\citet{goyal2008comprehensive} famously showed that most economic variables
studied over~30 years could not reliably predict the stock market out-of-sample.
\citet{campbell2008predicting} pushed back, arguing that even modest
predictability is economically valuable and introducing the out-of-sample
$R^2_{OOS}$ metric. \citet{huang2022scaled} proposed ``Scaled PCA,'' which
weights features by predictive power before dimension reduction - a key
inspiration. \citet{zou2005regularization} developed the Elastic Net,
combining L1 (lasso) and L2 (ridge) penalties, well-suited for datasets with
many correlated predictors. Our work differs by combining screening, scaling,
factor extraction, and regime interaction into a single defensible pipeline,
and by emphasizing statistical rigor (bootstrap confidence intervals,
permutation tests, Diebold-Mariano tests) to avoid overclaiming on a noisy
problem.


\section{Dataset and Features}

\paragraph{Economic indicators.}
Our primary data source is the FRED-MD database \citep{mccracken2016fredmd},
a collection of~134 monthly macroeconomic time series maintained by the
Federal Reserve Bank of St.\ Louis (``FRED'' = Federal Reserve Economic Data).
The indicators span output (GDP, industrial production), labor (employment,
housing starts), consumption (retail sales), money/credit (M2, bank credit),
interest rates (Treasury yields, corporate spreads), prices (CPI, PPI),
and sentiment (consumer confidence, VIX). Data covers January~1979 through
June~2026. We fill missing values forward and cap extreme outliers at the~1st
and~99th percentiles.

\paragraph{Target variable.}
The target is the monthly percentage change of the S\&P~500 index (Yahoo
Finance, ticker \texttt{\^{}GSPC}): $\frac{P_t - P_{t-1}}{P_{t-1}} \times 100$.
Features are aligned so that economic data from month~$t$ predicts returns
for month~$t+1$, mimicking a realistic trading scenario where macro reports
are released with a lag.

\paragraph{Train/test splits.}
We use an expanding ``walk-forward'' design: at each step the model trains on
all data up to the current month and predicts the next month. The first~60
months form the initial training set. For a cleaner evaluation, we also fix a
split at January~2000 - the model trains only on pre-2000 data and is
evaluated strictly on 2000--2026 with no retraining (``true OOS'').


\section{Methods}

The SSRF pipeline has four stages, each addressing a specific challenge in
macroeconomic forecasting:

\paragraph{Stage 1: Group-wise Screening.}
We have~134 features but only~540 months of data. Features are grouped into
economic categories (output, labor, inflation, etc.), and within each
category we retain only those with a t-statistic magnitude above~0.75. This
prevents any single category from dominating. About~40--60 features survive.

\paragraph{Stage 2: Predictive Scaling (``Scaled PCA'').}
Before extracting factors, each surviving feature is scaled by its univariate
regression slope against returns. This step, inspired by
\citet{huang2022scaled}, amplifies predictive features and suppresses noise.

\paragraph{Stage 3: Factor Extraction (PCA).}
Even after screening, the remaining features are highly correlated (they all
track the same economy). Principal Component Analysis (PCA) extracts~10
summary factors capturing the main independent sources of variation. The first
factor typically tracks the business cycle; later factors capture interest-rate
or inflation dynamics.

\paragraph{Stage 4: Regime Interaction.}
Markets behave differently in calm vs.\ turbulent periods. We measure the
current ``regime'' using rolling~12-month volatility - essentially asking
``is this market calmer or more volatile than usual?'' Each factor is
multiplied by this regime measure, creating interaction terms that let the
model learn different behavior under different conditions.

\paragraph{Final Regression.}
The regressor is Elastic Net \citep{zou2005regularization}, which combines L1
(lasso, driving irrelevant coefficients to zero) and L2 (ridge, shrinking
coefficients without eliminating them). We use $\alpha=0.001$ and
$\ell_1$-ratio$=0.5$ (equal mix). Training is walk-forward: the window
expands with each month.

\paragraph{Baselines.}
We compare against historical average (always predict the training mean),
momentum (predict the same direction as last month), and random guessing.

\paragraph{Asymmetric leverage.}
We also test asymmetric position sizing: 2.5$\times$ long when predicting
positive returns but only 0.25$\times$ short when predicting negative returns.
This exploits the empirical observation that long signals are more reliable
than short signals in macro forecasting.


\section{Experiments / Results / Discussion}

\paragraph{Evaluation metrics.}
We use four primary metrics:
\begin{itemize}
  \item \textbf{Hit ratio (direction accuracy):} percentage of months where the
    predicted sign matches the actual sign (random =~50\%).
  \item \textbf{Sharpe ratio:} annualized mean return divided by annualized
    volatility, measuring risk-adjusted return.
  \item \textbf{$R^2_{OOS}$:} compares the model's squared errors against the
    historical average (positive = beats the average).
  \item \textbf{Bootstrap 95\% confidence interval for Sharpe:} computed by
    resampling the return stream 1000 times (excludes zero if significant).
\end{itemize}

\paragraph{Main results.}
Table~\ref{tab:main} shows the primary out-of-sample results (2000--2026).

\begin{table}[ht]
\centering
\caption{Out-of-sample results (2000--2026).}
\label{tab:main}
\begin{tabular}{lcccc}
\toprule
Model & Hit Ratio & Sharpe & $R^2_{OOS}$ & Total Return \\
\midrule
SSRF (full)   & \textbf{70.3\%} & \textbf{0.51} & $-$0.005 & $+$3,808\% \\
Momentum      & 62.7\%          & 0.50          & $-$0.003 & $+$30\%   \\
Always Long   & 62.5\%          & ---           & 0.000    & ---       \\
S\&P~500 B\&H & ---             & 0.66          & ---      & $+$444\%  \\
\bottomrule
\end{tabular}
\end{table}

SSRF achieves~70.3\% direction accuracy - correctly predicting up~vs~down in
~7 of every~10 months. The first baseline worth beating is not~50\% (random
coin-flip) but~63\%: over this period the S\&P~500 went up in~198 of~317
months, so an ``always long'' strategy is right~62.5\% of the time simply by
riding the secular bull market. Momentum and the rolling historical mean hit
the same~63\% as this naive baseline. SSRF's~70.3\% is~8 percentage points
above the always-long benchmark, meaning it genuinely adds predictive value
rather than just betting on market drift. The Sharpe ratio is~0.51 (95\% CI:
[0.09,~1.21]), statistically significant ($p<0.0001$). The $R^2_{OOS}$ is
slightly negative ($-$0.005), meaning magnitude accuracy is comparable to the
historical average. This is typical in return prediction - direction is
easier to predict than exact magnitude, but directional bets still generate
substantial cumulative profits.

Walk-forward configurations all beat baselines: expanding window gives~66\%
hit ratio (Sharpe~0.60); fixed~60-month window gives~63\% (Sharpe~0.53).

\paragraph{Asymmetric leverage.}
Table~\ref{tab:leverage} shows the effect of 2.5$\times$ long~/~0.25$\times$
short leverage with enhanced features including the CAPE ratio and put/call
ratio:

\begin{table}[ht]
\centering
\caption{Asymmetric leverage (2.5$\times$/0.25$\times$).}
\label{tab:leverage}
\begin{tabular}{lcccc}
\toprule
Model & Ann.\ Return & Ann.\ Vol. & Sharpe & Max.\ Drawdown \\
\midrule
XGBoost & 8.6\% & 10.4\% & \textbf{0.84} & $-$9.1\% \\
Ensemble & 10.1\% & 12.8\% & 0.81 & $-$11.5\% \\
Linear & 9.4\% & 13.1\% & 0.75 & $-$9.4\% \\
S\&P~500 B\&H & 8.9\% & 13.4\% & 0.66 & $-$46.7\% \\
\bottomrule
\end{tabular}
\end{table}

Asymmetric leverage boosts Sharpe substantially (e.g., XGBoost:~0.84 vs.\
$<$0.7 symmetric) while cutting maximum drawdown from~$-$47\% to~$-$9\%.

When applied to individual sectors, the model performs best on defensive
sectors: Consumer Staples ($R^2_{OOS}=+0.04$,~64\% hit ratio) and Healthcare
($R^2_{OOS}=+0.03$). It struggles with cyclical sectors like Energy
($R^2_{OOS}=-1.94$) and Financials ($R^2_{OOS}=-1.04$).

\paragraph{What didn't work.}
Several experiments failed or disappointed. The regime interaction stage
(Stage~4) was tested early and disabled for final experiments because it
consistently hurt out-of-sample performance - the interaction terms added
more noise than signal. Extreme leverage (25$\times$ long) destroyed every
portfolio, with 72--84\% drawdowns across all model types. The $R^2_{OOS}$
remains negative in all configurations - the model cannot predict return
magnitudes, only direction.  Transaction costs, even at professional-tier
rates (15~bps), consume~1--2\% of annual returns, narrowing the advantage
over buy-and-hold. These failures reinforce a core lesson:
equity return prediction is a low signal-to-noise problem, and claimed
successes should always be viewed skeptically without proper out-of-sample
testing and statistical significance.


\section{Code Usage and Extensibility}

The codebase is designed for flexibility beyond the experiments described
above. It provides a command-line interface (CLI) and Python API for
exploring different model types, data sources, and trading configurations.
Table~\ref{tab:config} summarizes the key knobs available.

\begin{table}[ht]
\centering
\caption{Key configurable features.}
\label{tab:config}
\small
\begin{tabular}{lll}
\toprule
\textbf{Feature} & \textbf{CLI Flag} & \textbf{Options / Examples} \\
\midrule
Model type & \texttt{--model-type} & elasticnet, linear, xgboost,\\
  &  & random\_forest, catboost, mlp, ensemble \\
Screening threshold & \texttt{--t-stat-threshold} & 0.5~(lenient), 0.75~(default), 2.0~(strict) \\
PCA factors & \texttt{--n-factors} & 5, 10~(default), 20 \\
Regime window & \texttt{--regime-window} & 6~(fast), 12~(default), 24~(conservative) \\
Long/short leverage & \texttt{--max-long/--max-short} & 1.0/1.0~(symmetric), 2.5/0.25~(sweet spot) \\
Training window & \texttt{--train-window} & 60~(5yr), 120~(10yr), 180~(15yr) \\
Rebalance frequency & \texttt{--step-size} & 1~(monthly), 3~(quarterly) \\
Conviction filter & \texttt{--conviction-*} & 1.0~(default), 2.0~(only strong signals) \\
Account tier (TC) & \texttt{--account-tier} & micro, standard, professional, institutional \\
Margin rate & \texttt{--margin-rate} & 0.05~(5\%), 0.10~(10\%) \\
Drawdown limit & \texttt{--drawdown-limit} & 0.25~(reduce leverage at~25\% loss) \\
L1 vs L2 balance & \texttt{--l1-ratio} & 0.0~(pure Ridge), 0.5~(mix), 1.0~(pure Lasso) \\
Regularization & \texttt{--alpha} & 0.0001~(light), 0.001~(default), 0.1~(heavy) \\
\bottomrule
\end{tabular}
\end{table}

\paragraph{Model types.}
Seven final regressors are available. Four are covered in CS229 - linear
regression \texttt{--unregularized}, Ridge (implied by \texttt{--l1-ratio~0.0}),
Lasso (\texttt{--l1-ratio~1.0}), and Elastic Net (default). Three more are
tree-based: XGBoost (gradient-boosted trees, often best for non-linear
patterns), Random Forest, and a neural-network-style MLP. An Ensemble mode
combines all models via averaging. Switching is as simple as changing
\texttt{--model-type}.

\paragraph{Data sources.}
The pipeline fetches~34 raw FRED indicators and automatically derives~24
additional features (yield curve slopes, credit spreads, money supply growth
rates, sentiment regimes). It also loads three alternative indicators: the
CAPE ratio (Shiller's cyclically-adjusted P/E, available from~1881),
the CBOE put/call ratio (options market sentiment, from~1995), and
FINRA margin debt (broker-dealer leverage, from~1945). A VIX volatility proxy
fills pre-1990 values using realized volatility. To add new indicators, one
simply appends FRED series IDs to the configuration's indicator groups.

\paragraph{Sector rotation.}
The \texttt{--sector-rotation} flag lets users predict any of~11 economic
sectors (Technology, Healthcare, Energy, Financials, etc.) instead of the
S\&P~500, producing~77 sector-specific features (relative momentum, RSI,
cross-sectional rank, correlation, and trend signals).

\paragraph{Quick-start commands.}
\begin{verbatim}
# Default SSRF
python -m src.main --sample-data

# Compare model types
python -m src.main --model-type xgboost

# Asymmetric leverage with transaction costs
python -m src.main --max-long 2.5 --max-short 0.25 \\
    --tc-backtest --account-tier professional

# Sector rotation with alternative features
python -m src.main --sector-rotation Technology

# Aggressive: short training, fast rebalance
python -m src.main --train-window 36 --step-size 1

# Conservative: strict screening + conviction filter
python -m src.main --t-stat-threshold 2.0 --conviction-filter \\
    --conviction-threshold 2.0

# Full test suite
python test_all.py

# 7-model comparison with leverage sweep
python run_all_models_oos.py --sweep

# Real data OOS test
python run_oos_real_data.py
\end{verbatim}


\section{Conclusion / Future Work}

We built a multi-stage pipeline that achieves~70\% directional accuracy and a
statistically significant Sharpe of~0.51 out-of-sample from 2000--2026.
Asymmetric leverage boosts the Sharpe to~0.84 while nearly eliminating large
drawdowns.

Three limitations suggest future work. First, we do not model the release lag
of economic data - FRED statistics for a given month are often published
weeks later. Second, our regime detection (simple volatility percentile) could
be replaced with a hidden Markov model for more accurate state estimation.
Third, transaction costs were tested only in a limited analysis; a production
system would need to account for slippage and market impact.

\vspace{2em}

\section{Contributions}
Sole author. All code, analysis, and writing by Jagat Brahma.

\bibliographystyle{plainnat}
\bibliography{references}

\end{document}
